\documentclass[10pt,a4paper]{article}

\usepackage[T1]{fontenc}
\usepackage[utf8]{inputenc}
\usepackage[spanish]{babel}
\usepackage{lmodern}
\usepackage{geometry}
\usepackage{graphicx}
\usepackage{booktabs}
\usepackage{tabularx}
\usepackage{array}
\usepackage{amsmath}
\usepackage{enumitem}
\usepackage{xcolor}
\usepackage{fancyhdr}
\usepackage{lastpage}
\usepackage{pdflscape}
\usepackage{pgfgantt}
\usepackage[hidelinks]{hyperref}

\geometry{top=1.65cm,bottom=1.55cm,left=1.65cm,right=1.65cm}
\setlength{\parindent}{0pt}
\setlength{\parskip}{0.35em}
\setlist{nosep,leftmargin=1.25em}
\renewcommand{\arraystretch}{1.15}

\definecolor{upcblue}{HTML}{0076A8}
\definecolor{deepblue}{HTML}{174A6E}
\definecolor{lightblue}{HTML}{EAF4F8}
\definecolor{softgray}{HTML}{F2F4F5}
\definecolor{amber}{HTML}{D98E04}
\definecolor{green}{HTML}{5D8A66}

\pagestyle{fancy}
\fancyhf{}
\lhead{Entrega 2 de GEP · Planificación temporal}
\rhead{Víctor Abelló Torres}
\cfoot{\thepage\ de \pageref{LastPage}}
\renewcommand{\headrulewidth}{0.35pt}

\newcommand{\task}[1]{\textbf{#1}}
\newcommand{\tbd}{\textbf{TBD}}
\newcolumntype{Y}{>{\raggedright\arraybackslash}X}
\newcolumntype{C}[1]{>{\centering\arraybackslash}p{#1}}
\newcolumntype{L}[1]{>{\raggedright\arraybackslash}p{#1}}

\hypersetup{
  pdftitle={Entrega 2 de GEP: planificación temporal},
  pdfauthor={Víctor Abelló Torres}
}

\begin{document}
\hypersetup{pageanchor=false}
\begin{titlepage}
  \thispagestyle{empty}
  \centering
  \includegraphics[width=0.30\textwidth]{../../latex/assets/upc-logo.pdf}\par
  \vspace{1.45cm}
  {\color{upcblue}\Large\bfseries Entrega 2 de GEP\par}
  \vspace{0.3cm}
  {\Huge\bfseries Planificación temporal\par}
  \vspace{1.15cm}
  {\Large Estimación nutricional de platos a partir de imágenes mediante modelos de aprendizaje profundo: evaluación para su ejecución en dispositivos edge\par}
  \vfill
  {\Large\bfseries Víctor Abelló Torres\par}
  \vspace{0.8cm}
  \begin{tabular}{rl}
    Titulación: & Grado en Ingeniería Informática (Computación) \\
    Director: & Sergi Ramírez Mitjans \\
    Codirector: & Dante Conti \\
    Tutora de GEP: & Yadira Alejandra Salazar \\
    Fecha de entrega: & 28/09/2026
  \end{tabular}
  \vfill
  Facultat d'Informàtica de Barcelona\\
  Universitat Politècnica de Catalunya -- BarcelonaTech
\end{titlepage}

\hypersetup{pageanchor=true}
\pagenumbering{arabic}

{\color{deepblue}\Large\bfseries 1. Marco de la planificación}

\vspace{0.25em}
La planificación cubre el proyecto completo, desde la definición del protocolo experimental hasta la preparación de la defensa. Se adopta un enfoque iterativo e incremental: primero se obtiene un producto mínimo viable con datos, línea base y evaluación; después se incorporan tres familias visuales; finalmente se comparan precisión y coste, se selecciona un modelo y se valida su inferencia en un entorno limitado. La documentación y el seguimiento se realizan de forma concurrente durante todo el proyecto.

\begin{table}[h]
\centering
\caption{Datos globales de la planificación.}
\small
\begin{tabularx}{\textwidth}{>{\bfseries}L{4.0cm}Y}
\toprule
Inicio planificado & 14/09/2026 \\
Fin técnico del proyecto & 15/01/2027 \\
Fin desde el punto de vista de la FIB & 29/01/2027 (fecha objetivo provisional) \\
Lectura prevista del TFG & 29/01/2027, pendiente de confirmación en Racó (\tbd) \\
Esfuerzo total & 540 horas = 90 jornadas equivalentes de 6 horas \\
Ritmo ordinario & 6 h/día, de lunes a viernes; 30 h/semana \\
Horizonte temporal & 20 semanas naturales; la diferencia entre capacidad y esfuerzo absorbe festivos, indisponibilidad y riesgo \\
Responsable & Estudiante; supervisión del director y codirector \\
\bottomrule
\end{tabularx}
\end{table}

\textbf{Criterio de estimación.} Las horas incluyen implementación, comprobación, registro y cierre de cada tarea. Los entrenamientos se estiman por tiempo de trabajo activo; el cómputo desatendido no se contabiliza como dedicación, aunque condiciona el calendario. Ninguna tarea técnica alcanza 50 horas; la memoria se mantiene como una única tarea transversal de 60 horas porque su actividad es homogénea y recurrente.

\textbf{Organización iterativa.} El trabajo se divide en iteraciones de dos semanas. Al final de cada una se revisan resultados, riesgos y trabajo pendiente en el tablero Kanban. El Gantt fija el marco global, mientras que el tablero permite reordenar tareas sin perder los hitos ni la fecha final.

\textbf{Recursos.} R1: estudiante; R2: director y codirector; R3: portátil, Python, TensorFlow/Keras y Git; R4: Nutrition5k, almacenamiento y entorno de cómputo/Colab; R5: LiteRT y dispositivo edge (modelo concreto, \tbd); R6: Notion, LaTeX y herramientas FIB.

\newpage
{\color{deepblue}\Large\bfseries 2. Tareas, estimaciones y recursos (I)}

\vspace{0.2em}
La tabla describe tareas verificables y no simples fases. La columna ``resultado'' concreta qué debe existir para considerar terminada cada actividad; ``Pred.'' identifica sus predecesoras.

\begin{table}[h]
\centering
\caption{Tareas T01--T12.}
\scriptsize
\begin{tabularx}{\textwidth}{C{0.9cm}Y C{0.65cm} L{1.5cm} L{1.9cm}}
\toprule
\textbf{Cód.} & \textbf{Tarea y resultado de cierre} & \textbf{h} & \textbf{Pred.} & \textbf{Recursos} \\
\midrule
T01 & \task{Fijar alcance y protocolo.} Objetivos, variables, métricas y criterios de comparación aprobados. & 14 & -- & R1, R2, R6 \\
T02 & \task{Revisar bibliografía dirigida.} Matriz de trabajos, modelos, datos, métricas y límites relevantes. & 24 & -- & R1, R3, R6 \\
T03 & \task{Preparar el entorno reproducible.} Repositorios, dependencias, pruebas y estructura de experimentos operativos. & 10 & T01 & R1, R3 \\
T04 & \task{Adquirir y validar Nutrition5k.} Descarga controlada, inventario, integridad y muestra inspeccionada. & 18 & T01 & R1, R3, R4 \\
T05 & \task{Construir el flujo de datos.} Extracción de fotogramas, carga, transformación y caché reproducibles. & 28 & T03, T04 & R1, R3, R4 \\
T06 & \task{Definir particiones y controles de fuga.} División fijada por plato y comprobaciones automatizadas. & 14 & T05 & R1, R3, R4 \\
T07 & \task{Implementar y evaluar la línea base.} Predictor medio, métricas y artefactos de referencia guardados. & 14 & T06 & R1, R3, R4 \\
T08 & \task{Implementar la CNN de regresión.} Arquitectura y cinco salidas integradas en el flujo común. & 24 & T06, T07 & R1, R3, R4 \\
T09 & \task{Entrenar y ajustar la CNN.} Configuración acotada, mejor punto de control y registro de resultados. & 34 & T08 & R1, R3, R4 \\
T10 & \task{Adaptar MobileNet.} Backbone eficiente, cabeza de regresión y estrategia de ajuste definidos. & 24 & T06, T07 & R1, R3, R4 \\
T11 & \task{Entrenar y ajustar MobileNet.} Experimentos comparables, punto de control y resultados registrados. & 34 & T10 & R1, R3, R4 \\
T12 & \task{Adaptar un transformador visual.} Modelo, cabeza de regresión y configuración compatibles con el protocolo. & 26 & T06, T07 & R1, R3, R4 \\
\bottomrule
\end{tabularx}
\end{table}

\textbf{Secuencia lógica inicial.} T01 permite configurar el entorno y los datos; por tanto, \(\text{T01}\rightarrow\{\text{T03,T04}\}\). La cadena \(\{\text{T03,T04}\}\rightarrow\text{T05}\rightarrow\text{T06}\rightarrow\text{T07}\) produce el MVP experimental. Después de T06 y T07, las tres familias se implementan en ramas de trabajo independientes: \(\text{T08}\rightarrow\text{T09}\), \(\text{T10}\rightarrow\text{T11}\) y \(\text{T12}\rightarrow\text{T13}\). Esta separación permite ejecutar unas tareas mientras otra familia está entrenando.

\newpage
{\color{deepblue}\Large\bfseries 3. Tareas, estimaciones y dependencias (II)}

\begin{table}[h]
\centering
\caption{Tareas T13--T23.}
\scriptsize
\begin{tabularx}{\textwidth}{C{0.9cm}Y C{0.65cm} L{1.65cm} L{1.85cm}}
\toprule
\textbf{Cód.} & \textbf{Tarea y resultado de cierre} & \textbf{h} & \textbf{Pred.} & \textbf{Recursos} \\
\midrule
T13 & \task{Entrenar y ajustar el transformador.} Experimentos comparables, punto de control y registro completo. & 36 & T12 & R1, R3, R4 \\
T14 & \task{Aplicar la evaluación común.} Predicciones, métricas por objetivo y análisis de errores incorporados al terminar cada modelo. & 28 & T07; cierre tras T09, T11, T13 & R1, R3, R4 \\
T15 & \task{Medir el coste de inferencia.} Tamaño, latencia y, cuando sea posible, memoria bajo condiciones documentadas. & 24 & T14 & R1, R3, R4 \\
T16 & \task{Convertir y cuantizar el candidato.} Modelo LiteRT válido y comparación antes/después de la conversión. & 24 & T15 & R1, R3, R5 \\
T17 & \task{Validar la ejecución edge.} Inferencia, latencia y limitaciones medidas en el entorno seleccionado. & 20 & T16 & R1, R3, R5 \\
T18 & \task{Sintetizar la comparación.} Tablas de precisión--coste, interpretación y respuesta a la pregunta de investigación. & 20 & T14, T15, T17 & R1, R3, R6 \\
T19 & \task{Analizar sostenibilidad y ética.} Riesgos de privacidad, sesgo, seguridad, energía y uso indebido. & 12 & T15 & R1, R2, R6 \\
T20 & \task{Redactar y revisar la memoria.} Texto, figuras y trazabilidad actualizados durante todo el proyecto. & 60 & T01; cierre tras T18, T19 & R1, R2, R6 \\
T21 & \task{Control y seguimiento.} Revisión semanal de horas, avance, riesgos, acuerdos y replanificación. & 24 & T01 & R1, R2, R6 \\
T22 & \task{Cerrar la entrega académica.} Revisión final, portada, anexos, depósito y trámites FIB aplicables. & 10 & T18, T19, T20 & R1, R2, R6 \\
T23 & \task{Preparar la defensa.} Diapositivas, demostración opcional, ensayos y reserva de aula si procede. & 18 & T22 & R1, R2, R3, R6 \\
\midrule
 & \textbf{Total} & \textbf{540} & & \\
\bottomrule
\end{tabularx}
\end{table}

\textbf{Dependencias de cierre.} Las mediciones se realizan cuando los modelos ya se han evaluado: \(\{\text{T09,T11,T13}\}\rightarrow\text{T14}\rightarrow\text{T15}\rightarrow\text{T16}\rightarrow\text{T17}\). La conclusión experimental requiere T14, T15 y T17. El cierre técnico reúne T18, T19 y la versión final de T20; después se realizan los trámites (T22) y la preparación de la defensa (T23).

\begin{table}[h]
\centering
\caption{Hitos de control.}
\small
\begin{tabularx}{\textwidth}{L{1.4cm}L{2.7cm}Y}
\toprule
\textbf{Hito} & \textbf{Fecha objetivo} & \textbf{Criterio observable} \\
\midrule
M1 & 02/10/2026 & MVP reproducible: datos, particiones, línea base y métricas. \\
M2 & 11/12/2026 & Tres familias entrenadas y evaluadas con el mismo protocolo. \\
M3 & 15/01/2027 & Comparación final, validación edge y memoria técnicamente cerrada. \\
M4 & 29/01/2027 & Depósito y defensa completados; fecha oficial pendiente de Racó. \\
\bottomrule
\end{tabularx}
\end{table}

\thispagestyle{fancy}
\clearpage
\begin{landscape}
\thispagestyle{fancy}
{\color{deepblue}\Large\bfseries 4. Diagrama de Gantt}

\vspace{0.2em}
Cada columna representa una semana. Las barras largas de documentación y seguimiento muestran la concurrencia; no implican dedicación diaria continua. La semana 20 corresponde al cierre institucional y la defensa provisional.

\centering
\begin{ganttchart}[
  hgrid,
  vgrid,
  x unit=0.79cm,
  y unit chart=0.46cm,
  y unit title=0.58cm,
  title height=1,
  bar height=0.60,
  group height=0.60,
  milestone height=0.55,
  title/.style={fill=deepblue,draw=white},
  title label font=\color{white}\bfseries\scriptsize,
  bar label font=\scriptsize,
  group label font=\scriptsize\bfseries,
  milestone label font=\scriptsize\bfseries,
  bar/.append style={fill=upcblue!72,draw=upcblue},
  group/.append style={fill=deepblue!82,draw=deepblue},
  milestone/.append style={fill=amber,draw=amber},
  bar label node/.append style={align=right,text width=5.4cm},
  group label node/.append style={align=right,text width=5.4cm},
  milestone label node/.append style={align=right,text width=5.4cm}
]{1}{20}
\gantttitle{Sep.}{3}\gantttitle{Oct.}{4}\gantttitle{Nov.}{4}\gantttitle{Dic.}{5}\gantttitle{Ene.}{4}\\
\gantttitlelist{1,...,20}{1}\\
\ganttgroup{Preparación y MVP}{1}{3}\\
\ganttbar{T01 Alcance y protocolo}{1}{2}\\
\ganttbar{T02 Revisión bibliográfica}{1}{5}\\
\ganttbar{T03 Entorno reproducible}{1}{1}\\
\ganttbar{T04 Adquisición y validación}{1}{2}\\
\ganttbar{T05 Flujo de datos}{1}{3}\\
\ganttbar{T06 Particiones y fugas}{3}{3}\\
\ganttbar{T07 Línea base}{3}{3}\\
\ganttmilestone{M1 MVP}{3}\\
\ganttgroup{Modelos y evaluación}{4}{17}\\
\ganttbar{T08 CNN: implementación}{4}{4}\\
\ganttbar{T09 CNN: entrenamiento}{5}{6}\\
\ganttbar{T10 MobileNet: adaptación}{7}{7}\\
\ganttbar{T11 MobileNet: entrenamiento}{8}{9}\\
\ganttbar{T12 Transformer: adaptación}{10}{10}\\
\ganttbar{T13 Transformer: entrenamiento}{11}{12}\\
\ganttbar{T14 Evaluación común}{6}{13}\\
\ganttmilestone{M2 Modelos comparados}{13}\\
\ganttbar{T15 Coste de inferencia}{14}{14}\\
\ganttbar{T16 Conversión y cuantización}{15}{16}\\
\ganttbar{T17 Validación edge}{17}{17}\\
\ganttbar{T18 Síntesis comparativa}{17}{18}\\
\ganttbar{T19 Sostenibilidad y ética}{15}{18}\\
\ganttgroup{Gestión y cierre}{1}{20}\\
\ganttbar[bar/.append style={fill=green!72,draw=green!85!black}]{T20 Memoria}{1}{19}\\
\ganttbar[bar/.append style={fill=green!72,draw=green!85!black}]{T21 Seguimiento}{1}{19}\\
\ganttmilestone{M3 Fin técnico}{18}\\
\ganttbar{T22 Entrega y trámites}{19}{19}\\
\ganttbar{T23 Defensa}{19}{20}\\
\ganttmilestone{M4 Fin FIB}{20}
\end{ganttchart}

\vfill
\raggedright\footnotesize
\textbf{Lectura del calendario.} S1 comienza el 14/09/2026; S3, el 28/09; S7, el 26/10; S11, el 23/11; S15, el 21/12; S17, el 04/01/2027; y S20, el 25/01. El trabajo activo se ajustará a 540 horas, por lo que festivos y esperas de cómputo no se sustituyen automáticamente por nuevas tareas.
\end{landscape}

\clearpage
{\color{deepblue}\Large\bfseries 5. Seguimiento, obstáculos y planes alternativos}

\vspace{0.2em}
Cada viernes se compararán horas reales y planificadas, tareas terminadas y riesgos. Una desviación superior al 15\,\% en horas o una semana en una tarea de la cadena de cierre activa una revisión con la dirección. Se reestima primero el trabajo pendiente; las tareas opcionales nunca desplazan el alcance mínimo.

\begin{table}[h]
\centering
\caption{Planes alternativos; no forman parte del Gantt base.}
\scriptsize
\begin{tabularx}{\textwidth}{L{2.25cm}Y L{3.0cm} L{2.35cm}}
\toprule
\textbf{Obstáculo / señal} & \textbf{Tareas alternativas} & \textbf{Efecto sobre duración} & \textbf{Recursos adicionales} \\
\midrule
\task{R1. Cómputo insuficiente.} Colas o entrenamientos que exceden una iteración. & \task{B1:} congelar más capas, limitar la búsqueda de hiperparámetros y entrenar con resolución o lote menores; mantener el mismo conjunto de prueba. & Consume hasta 18 h del margen; no cambia M3. & Acelerador externo; coste, si existe, \tbd. \\
\task{R2. Datos o almacenamiento.} Descarga incompleta, corrupción o falta de espacio. & \task{B2:} continuar con el piloto y artefactos cacheados mientras se recupera la copia completa; verificar hashes e inventario antes de reanudar. & Hasta 12 h de margen; T09/T11/T13 no cierran con el piloto. & Almacenamiento temporal R4. \\
\task{R3. Entorno edge no disponible o incompatible.} T16 no genera un modelo ejecutable. & \task{B3:} probar operadores alternativos, seleccionar el siguiente modelo compatible y medir LiteRT en CPU representativa; documentar que no equivale al dispositivo final. & Hasta 16 h; puede reducir el alcance de T17, no la comparación principal. & Equipo alternativo o préstamo, \tbd. \\
\task{R4. Error de los modelos superior al esperado.} Ninguna familia mejora claramente la línea base. & \task{B4:} auditar etiquetas, escalas y pérdidas; ejecutar un ajuste acotado. Si persiste, analizar el resultado negativo sin ocultarlo. & Hasta 14 h; no altera la fecha final. & R1, R2 y registros R4. \\
\task{R5. Retraso acumulado.} Se agota el margen o dos hitos consecutivos se incumplen. & \task{B5:} eliminar el demostrador y cualquier extensión multimodal; conservar línea base, CNN, MobileNet, transformer, evaluación común y un candidato edge. & Recupera hasta 60 h de trabajo opcional y protege M3/M4. & No requiere recursos nuevos. \\
\bottomrule
\end{tabularx}
\end{table}

\textbf{Reserva y reducción de alcance.} El calendario ofrece aproximadamente 60 horas de capacidad no asignada respecto al ritmo ordinario de 20 semanas. No se presenta como trabajo adicional ni como aumento del presupuesto de 540 horas: absorbe festivos, esperas y activaciones de B1--B4. Si la desviación supera esa reserva, se aplica B5. La profundidad, el texto, el demostrador visual y una búsqueda extensa de hiperparámetros siguen siendo opcionales.

\textbf{Dos finales explícitos.} M3 marca el final técnico: la pregunta de investigación queda respondida con resultados trazables. M4 marca el final académico: depósito, trámites y defensa. Si Racó asigna otra fecha, T22 y T23 se desplazarán manteniendo M3 y el trabajo técnico ya cerrado.

\textbf{Control de calidad.} Cada hito requiere artefactos reproducibles, una revisión de coherencia entre memoria y resultados y la actualización del registro de decisiones. Las reuniones dejan acuerdos y acciones; las horas reales alimentan la siguiente replanificación.

\end{document}
